\section{基于反步的全状态反馈}

本节记录一种用于强制
\[
\dot{\bm p}_{i0}\to -\kappa_i \bm p_{i0}+\bm\phi_i,
\]
的反步式尝试，并解释为何该构造对当前的异构高阶模型不闭合。

\subsection{理想的二阶尝试}

定义第一个反步误差为
\[
\bm e_{1,i}:=\dot{\bm p}_{i0}+\kappa_i\bm p_{i0}-\bm\phi_i.
\]
若 $\bm e_{1,i}\to\bm 0$，则确实有
\[
\dot{\bm p}_{i0}\to -\kappa_i\bm p_{i0}+\bm\phi_i.
\]

对于理想化的二阶相对动力学
\[
\ddot{\bm p}_{i0}=\bm u_i+\bm d_i-\ddot{\bm p}_0,
\]
误差动力学为
\[
\dot{\bm e}_{1,i}
=\bm u_i+\bm d_i-\ddot{\bm p}_0
 +\kappa_i\dot{\bm p}_{i0}-\dot{\bm\phi}_i.
\]
若 $\dot{\bm\phi}_i$ 和 $\ddot{\bm p}_0$ 可得，则理想的反步选择为
\[
\bm u_i
=-\kappa_i\dot{\bm p}_{i0}+\dot{\bm\phi}_i+\ddot{\bm p}_0
-k_i\bm e_{1,i}-\rho_i\operatorname{sgn}(\bm e_{1,i}),
\]
其中 $k_i>0$ 且 $\rho_i>\overline d_i$。
则
\[
\dot{\bm e}_{1,i}
=-k_i\bm e_{1,i}+\bm d_i-\rho_i\operatorname{sgn}(\bm e_{1,i}),
\]
李雅普诺夫函数 $V_i=\frac12\|\bm e_{1,i}\|^2$ 满足
\[
\dot V_i
\le -k_i\|\bm e_{1,i}\|^2-(\rho_i-\overline d_i)\|\bm e_{1,i}\|.
\]
因此 $\bm e_{1,i}\to\bm 0$，且在理想化设置中实现了期望的阶一跟踪律。

\subsection{为何该构造不闭合}

当前模型并非纯二阶系统。它是
\[
\bm p_i^{(m_i)}=\bm u_i+\bm d_i,
\qquad m_i\ge 1,
\]
因此完整的反步设计需要递归的虚拟控制链。
仅在第一步，上述设计就需要 $\dot{\bm\phi}_i$；
对于 $m_i>2$，递归需要
\[
\bm\phi_i,\ \dot{\bm\phi}_i,\ \ddot{\bm\phi}_i,\ \dots,\ \bm\phi_i^{(m_i-1)}.
\]
这正是尝试变得脆弱之处：

\begin{itemize}
    \item $\bm\phi_i$ 依赖于时变邻域集合 $\mathbb N_i(t)$，
          因此它仅分段光滑。
    \item 在 $\|\bm p_{ij}\|=d$ 附近，排斥项变得奇异，
          其导数可以增长得非常快。
    \item 所需的目标导数 Beyond $\ddot{\bm p}_0$ 在模型中并未假设可得。
    \item 分布式实现必须评估所有成对排斥项的高阶导数，
          这繁琐且数值上脆弱。
\end{itemize}

因此，上述反步控制器仅是对理想二阶情形有用的设计思路。
它并未为本文考虑的异构高阶系统给出一个简洁的显式控制器。

\subsection{实际说明}

若将 $\bm\phi_i$ 替换为滤波信号 $\bm\Phi_{1,i}$，
导数奇异性被软化，但此时设计退化为先前的滤波 APF 框架。
在这种情况下，反步层可稳定
\[
\dot{\bm p}_{i0}\approx -\kappa_i\bm p_{i0}+\bm\Phi_{1,i},
\]
然而除非引入额外的不变性论证，否则直接的障碍能量证明仍会丢失。

\subsection{可能的替代方案}

失败的精确反步尝试仍有用，因为它识别了主要障碍：
反复对原始 APF 项求导不是一种稳健的操作。
若干替代路线可避免此障碍。

\paragraph{指令滤波反步。}
将期望的阶一行为设为
\[
\bm v_{i,d}:=-\kappa_i\bm p_{i0}+\bm\phi_i,
\]
但不解析地微分 $\bm v_{i,d}$。
相反，将 $\bm v_{i,d}$ 通过稳定的指令滤波器，并在反步递归中
使用滤波器输出及其可用的导数。
这接近动态面控制。
它减少了对 $\dot{\bm\phi}_i,\dots,\bm\phi_i^{(m_i-1)}$ 的需求，
因此更兼容邻域切换和奇异排斥场。
代价是闭环跟踪的是 $\bm v_{i,d}$ 的滤波版本，
因此需要额外的滤波误差分析。

\paragraph{外环/内环设计。}
使用
\[
\dot{\bm p}_{i0}= -\kappa_i\bm p_{i0}+\bm\phi_i
\]
仅作为外环参考模型。
内环随后使用高增益、滑模或反步控制器使 $\dot{\bm p}_{i0}$ 尽可能
精确地跟踪该参考。
这给出了清晰的模块化解释：
APF 产生期望的阶一围捕运动，而高阶动力学由内环处理。
代价是除非内环跟踪误差被精确驱动到零，否则实现的运动是近似的。

\paragraph{基于一阶误差的高阶滑模。}
可直接定义
\[
\bm s_i:=\dot{\bm p}_{i0}+\kappa_i\bm p_{i0}-\bm\phi_i
\]
并使用高阶滑模律强制 $\bm s_i=\bm 0$。
这避免了构建完整的递归反步链。
然而，控制仍依赖于 $\bm\phi_i$ 的正则性假设以及从 $\bm u_i$ 到
$\bm s_i$ 的相对阶。
对于异构的 $m_i$，所需高阶滑模律随智能体阶数变化，且分析变为
依赖具体情形的。

\paragraph{控制障碍函数安全层。}
一条不同的路线是将收敛性与安全性分离。
反步或滤波 APF 控制器用作标称控制器，
而 CBF 或高阶 CBF 层强制
\[
h_{ij}(\bm p)=\|\bm p_{ij}\|^2-d^2>0.
\]
这直接针对安全集合的前向不变性，避免要求 APF 跟踪证明也证明避碰。
在所有替代方案中，这是获得严格全时安全证书的最简洁方式，
但它引入了一个额外的约束执行层。

\subsection{试探性方向}

最有希望的方向似乎是指令滤波反步与 CBF 安全层的结合。
指令滤波将使期望的阶一行为
\[
\dot{\bm p}_{i0}\approx -\kappa_i\bm p_{i0}+\bm\phi_i
\]
对异构高阶智能体可实现，而无需反复微分 $\bm\phi_i$。
CBF 层随后将提供缺失的避碰前向不变性证明。
这种组合保留了 APF 围捕的直觉，同时避免了精确反步尝试的
最弱点。